% GENERATED FILE. Edit canonical lessons, then run npm run generate:books.
% canonical-source-hash: fnv1a64:7ce580dcf1cc6789
% canonical-lessons: MR-C382-sambhavya, MR-C382-bahudha, MR-C382-sahjikac, MR-C382-niscit, MR-C382-bhasne

\chapter{Likely or Certain}
\label{ch:mr-likely-or-certain}
\hlchaptermodality{382}

\begin{chapteropening}
\textbf{By the end of this chapter:} \emph{I can say how likely or certain something is.}

Complete the last lesson of chapter 382: I can say how likely or certain something is.
\end{chapteropening}

\section[\texorpdfstring{\mr{संभाव्य} (sambhāvya)}{संभाव्य (sambhāvya)}]{\mr{संभाव्य} (sambhāvya) — likely; potential}

\label{lesson:MR-C382-sambhavya}



\begin{quote}
Before the new one: say the Marathi for so that, then the Marathi for as if.
\end{quote}

\subsection*{You'll want to know: \mr{संभाव्य}}
\textbf{\mr{संभाव्य}} — \emph{sambhāvya} — \textquotedblleft{}likely; potential\textquotedblright{}.

\begin{quote}
\mr{संभाव्य} \mr{धोका} \mr{लक्षात} \mr{घेऊन} \mr{शाळा} \mr{बंद} \mr{ठेवल्या}.
\end{quote}

\begin{quote}
\emph{sambhāvya dhokā lakṣāt gheūn śāḷā band ṭhevalyā.}
\end{quote}

\begin{quote}
Keeping the possible danger in mind, the schools were kept closed.
\end{quote}

Invariable; \mr{संभाव्यता}, probability.

\begin{cousinweb}
\textbf{\mr{संभाव्य}} (\emph{sambhāvya}) is a tatsama, taken from Sanskrit \emph{sambhāvya}, \textquotedblleft{}possible, probable\textquotedblright{} (Monier-Williams).
\end{cousinweb}

\begin{grammarlens}[title={What you've built}]
One more word for saying what happens if.
\end{grammarlens}

\subsection*{Your turn}
\begin{itemize}
\raggedright
  \item \emph{Say it:} \emph{sambhāvya}
  \item \emph{Say it:} \emph{sambhāvya}, once more
  \item \emph{Say it:} \emph{sambhāvya dhokā lakṣāt gheūn śāḷā band ṭhevalyā.} — Keeping the possible danger in mind, the schools were kept closed.
  \item \emph{Recall:} say the Marathi for so that, then the Marathi for as if, then say \emph{sambhāvya} again
\end{itemize}

\begin{tcolorbox}[breakable,colback=teal!4,colframe=teal!35!black,title={Before you move on}]
What does \mr{संभाव्य} mean? (\textquotedblleft{}Likely; potential\textquotedblright{}.) Say it once more.
\end{tcolorbox}
\section[\texorpdfstring{\mr{बहुधा} (bahudhā)}{बहुधा (bahudhā)}]{\mr{बहुधा} (bahudhā) — probably; most likely}

\label{lesson:MR-C382-bahudha}



\begin{quote}
Before the new one: say the Marathi for as if, then the Marathi for likely.
\end{quote}

\subsection*{You'll want to know: \mr{बहुधा}}
\textbf{\mr{बहुधा}} — \emph{bahudhā} — \textquotedblleft{}probably; most likely\textquotedblright{}.

\begin{quote}
\mr{बहुधा} \mr{उद्या} \mr{पाऊस} \mr{पडेल}.
\end{quote}

\begin{quote}
\emph{bahudhā udyā pāūs paḍel.}
\end{quote}

\begin{quote}
It will probably rain tomorrow.
\end{quote}

Often first in the sentence, before the guess: \mr{बहुधा} \mr{तो} \mr{येईल}.

\begin{cousinweb}
\textbf{\mr{बहुधा}} (\emph{bahudhā}) is a tatsama, taken from Sanskrit \emph{bahudhā}, \textquotedblleft{}in many ways\textquotedblright{} (CDIAL 9191); Marathi uses it for what is most likely.
\end{cousinweb}

\begin{grammarlens}[title={What you've built}]
One more word for saying what happens if.
\end{grammarlens}

\subsection*{Your turn}
\begin{itemize}
\raggedright
  \item \emph{Say it:} \emph{bahudhā}
  \item \emph{Say it:} \emph{bahudhā}, once more
  \item \emph{Say it:} \emph{bahudhā udyā pāūs paḍel.} — It will probably rain tomorrow.
  \item \emph{Recall:} say the Marathi for as if, then the Marathi for likely, then say \emph{bahudhā} again
\end{itemize}

\begin{tcolorbox}[breakable,colback=teal!4,colframe=teal!35!black,title={Before you move on}]
What does \mr{बहुधा} mean? (\textquotedblleft{}Probably; most likely\textquotedblright{}.) Say it once more.
\end{tcolorbox}
\section[\texorpdfstring{\mr{साहजिकच} (sāhjikac)}{साहजिकच (sāhjikac)}]{\mr{साहजिकच} (sāhjikac) — naturally; of course}

\label{lesson:MR-C382-sahjikac}



\begin{quote}
Before the new one: say the Marathi for likely, then the Marathi for probably.
\end{quote}

\subsection*{You'll want to know: \mr{साहजिकच}}
\textbf{\mr{साहजिकच}} — \emph{sāhjikac} — \textquotedblleft{}naturally; of course\textquotedblright{}.

\begin{quote}
\mr{तो} \mr{खूप} \mr{थकला} \mr{होता}, \mr{साहजिकच} \mr{लगेच} \mr{झोपला}.
\end{quote}

\begin{quote}
\emph{to khūp thaklā hotā, sāhjikac lagec jhoplā.}
\end{quote}

\begin{quote}
He was very tired, so naturally he fell asleep at once.
\end{quote}

\mr{साहजिक}, natural, + -\mr{च}.

\begin{cousinweb}
\textbf{\mr{साहजिकच}} (\emph{sāhjikac}) is built from \mr{साहजिक}, \textquotedblleft{}natural\textquotedblright{}, made on Sanskrit \emph{sahaja}, \textquotedblleft{}born together, innate\textquotedblright{} (Monier-Williams), with the stressing ending -\mr{च}.
\end{cousinweb}

\begin{grammarlens}[title={What you've built}]
One more word for saying what happens if.
\end{grammarlens}

\subsection*{Your turn}
\begin{itemize}
\raggedright
  \item \emph{Say it:} \emph{sāhjikac}
  \item \emph{Say it:} \emph{sāhjikac}, once more
  \item \emph{Say it:} \emph{to khūp thaklā hotā, sāhjikac lagec jhoplā.} — He was very tired, so naturally he fell asleep at once.
  \item \emph{Recall:} say the Marathi for likely, then the Marathi for probably, then say \emph{sāhjikac} again
\end{itemize}

\begin{tcolorbox}[breakable,colback=teal!4,colframe=teal!35!black,title={Before you move on}]
What does \mr{साहजिकच} mean? (\textquotedblleft{}Naturally; of course\textquotedblright{}.) Say it once more.
\end{tcolorbox}
\section[\texorpdfstring{\mr{निश्चित} (niścit)}{निश्चित (niścit)}]{\mr{निश्चित} (niścit) — certain; fixed}

\label{lesson:MR-C382-niscit}



\begin{quote}
Before the new one: say the Marathi for probably, then the Marathi for naturally.
\end{quote}

\subsection*{You'll want to know: \mr{निश्चित}}
\textbf{\mr{निश्चित}} — \emph{niścit} — \textquotedblleft{}certain; fixed\textquotedblright{}.

\begin{quote}
\mr{लग्नाची} \mr{तारीख} \mr{अजून} \mr{निश्चित} \mr{नाही}.
\end{quote}

\begin{quote}
\emph{lagnācī tārīkh ajūn niścit nāhī.}
\end{quote}

\begin{quote}
The wedding date is not fixed yet.
\end{quote}

\mr{निश्चित} \mr{करणे}, to fix; \mr{निश्चितपणे}, certainly.

\begin{cousinweb}
\textbf{\mr{निश्चित}} (\emph{niścit}) is a tatsama, taken from Sanskrit \emph{niścita}, \textquotedblleft{}ascertained, settled\textquotedblright{} (Monier-Williams).
\end{cousinweb}

\begin{grammarlens}[title={What you've built}]
One more word for saying what happens if.
\end{grammarlens}

\subsection*{Your turn}
\begin{itemize}
\raggedright
  \item \emph{Say it:} \emph{niścit}
  \item \emph{Say it:} \emph{niścit}, once more
  \item \emph{Say it:} \emph{lagnācī tārīkh ajūn niścit nāhī.} — The wedding date is not fixed yet.
  \item \emph{Recall:} say the Marathi for probably, then the Marathi for naturally, then say \emph{niścit} again
\end{itemize}

\begin{tcolorbox}[breakable,colback=teal!4,colframe=teal!35!black,title={Before you move on}]
What does \mr{निश्चित} mean? (\textquotedblleft{}Certain; fixed\textquotedblright{}.) Say it once more.
\end{tcolorbox}
\section[\texorpdfstring{\mr{भासणे} (bhāsṇe)}{भासणे (bhāsṇe)}]{\mr{भासणे} (bhāsṇe) — to seem; to appear}

\label{lesson:MR-C382-bhasne}



\begin{quote}
Before the new one: say the Marathi for naturally, then the Marathi for certain.
\end{quote}

\subsection*{You'll want to know: \mr{भासणे}}
\textbf{\mr{भासणे}} — \emph{bhāsṇe} — \textquotedblleft{}to seem; to appear\textquotedblright{}.

\begin{quote}
\mr{दुरून} \mr{ती} \mr{टेकडी} \mr{खूप} \mr{जवळ} \mr{भासते}.
\end{quote}

\begin{quote}
\emph{durūn tī ṭekḍī khūp javaḷ bhāste.}
\end{quote}

\begin{quote}
From afar that hill seems very close.
\end{quote}

Past (to him …): \mr{त्याला} \mr{भासलं}, \emph{tyālā bhāslaṁ}.

Formal; in speech \mr{वाटणे}. Also with the dative: \mr{मला} \mr{तसं} \mr{भासलं}.

\begin{cousinweb}
\textbf{\mr{भासणे}} (\emph{bhāsṇe}) is a tadbhava: CDIAL lists Marathi \emph{bhāsṇẽ}, \textquotedblleft{}to appear, seem, seem probable\textquotedblright{}, under Sanskrit \emph{bhāsati}, \textquotedblleft{}is bright, appears\textquotedblright{} (CDIAL 9481).
\end{cousinweb}

\begin{grammarlens}[title={What you've built}]
One more word for saying what happens if.
\end{grammarlens}

\subsection*{Your turn}
\begin{itemize}
\raggedright
  \item \emph{Say it:} \emph{bhāsṇe}
  \item \emph{Say it:} \emph{bhāsṇe}, once more
  \item \emph{Say it:} \emph{durūn tī ṭekḍī khūp javaḷ bhāste.} — From afar that hill seems very close.
  \item \emph{Recall:} say the Marathi for naturally, then the Marathi for certain, then say \emph{bhāsṇe} again
\end{itemize}

\begin{tcolorbox}[breakable,colback=teal!4,colframe=teal!35!black,title={Before you move on}]
What does \mr{भासणे} mean? (\textquotedblleft{}To seem; to appear\textquotedblright{}.) Say it once more.
\end{tcolorbox}
